% Bagian: first-order-logic
% Bab: models-theories
% Subbagian: size-of-structures

\documentclass[../../../include/open-logic-section]{subfiles}

\begin{document}

\olfileid[id]{fol}{mat}{siz}

\olsection{Mengungkapkan Ukuran \printtoken{P}{structure}}

\begin{explain}
Ada beberapa sifat struktur yang dapat kita ungkapkan bahkan tanpa
menggunakan simbol nonlogis suatu bahasa. Misalnya, ada !!{sentence}s yang
benar dalam !!a{structure} jika dan hanya jika !!{domain} dari !!{structure}
tersebut mempunyai sekurang-kurangnya, sebanyak-banyaknya, atau tepat sejumlah
tertentu~$n$ !!{element}s.
\end{explain}

\begin{prop}
!!^{sentence}
\begin{multline*}
  !A_{\ge n} \ident \lexists[x_1][\lexists[x_2][\dots\lexists[x_n][{}]]]\\
  \begin{aligned}
  (\eq/[x_1][x_2] \land {}
  \eq/[x_1][x_3] \land \eq/[x_1][x_4] \land \dots \land \eq/[x_1][x_n] \land {}\\
\eq/[x_2][x_3] \land \eq/[x_2][x_4] \land \dots \land {} \eq/[x_2][x_n] \land {} \\
\vdots\\
\eq/[x_{n-1}][x_n])
\end{aligned}
\end{multline*}
benar dalam !!a{structure}~$\Struct M$ jika dan hanya jika $\Domain M$
mempunyai sekurang-kurangnya $n$ !!{element}s. Akibatnya,
$\Sat{M}{\lnot !A_{\ge n+1}}$ jika dan hanya jika $\Domain M$ mempunyai
sebanyak-banyaknya~$n$ !!{element}s.

\end{prop}

\begin{prop}
!!^{sentence}
\begin{multline*}
!A_{= n} \ident \lexists[x_1][\lexists[x_2][\dots\lexists[x_n][{}]]] \\
\begin{aligned}
  (\eq/[x_1][x_2] \land {}
  \eq/[x_1][x_3] \land \eq/[x_1][x_4] \land \dots \land \eq/[x_1][x_n] \land {}\\
\eq/[x_2][x_3] \land \eq/[x_2][x_4] \land \dots \land {} \eq/[x_2][x_n] \land {} \\
\vdots\\
\eq/[x_{n-1}][x_n] \land {} \\
\lforall[y][(\eq[y][x_1] \lor \dots \lor \eq[y][x_n]]))
\end{aligned}
\end{multline*}
benar dalam !!a{structure}~$\Struct M$ jika dan hanya jika $\Domain M$
mempunyai tepat $n$ !!{element}s.
\end{prop}

\begin{prop}
Suatu !!{structure} tak hingga jika dan hanya jika struktur itu merupakan
model dari
\[
\{!A_{\ge 1}, !A_{\ge 2}, !A_{\ge 3}, \dots \}.
\]
\end{prop}

Tidak ada satu kalimat logis murni yang benar dalam~$\Struct M$ jika dan hanya
jika $\Domain M$ tak hingga. Namun, orang dapat memberikan !!{sentence}s
dengan !!{predicate}s nonlogis yang hanya mempunyai model-model tak hingga
(meskipun tidak setiap !!{structure} tak hingga merupakan modelnya). Sifat
sebagai struktur hingga dan sifat sebagai struktur !!{nonenumerable} bahkan
tidak dapat diungkapkan dengan suatu himpunan !!{sentence}s tak hingga.
Fakta-fakta ini mengikuti dari teorema kekompakan dan
L\"owenheim--Skolem.

\end{document}
